% !TEX program = xelatex
% 编译方式：在本目录执行两遍 xelatex（详见 docs/_manual_common/README.md）。
% 源码依据：src/scenario_generation/ 与 include/hacdcpf/analysis/（scenario_generation.hpp、
%           typhoon_resilience.hpp、typhoon_traffic_impact.hpp）源码及 docs/scenario_generation_contract.md，公式与参数以代码为准。
\documentclass[UTF8,fontset=fandol,a4paper,11pt]{ctexart}
\usepackage{hysim_manual}

\title{\textbf{交直流混合高弹性能源电力系统仿真平台}\\[0.5em]
       {\Large 场景生成技术手册}\\[0.3em]
       {\large 数学模型、接口参数、验证校核与工业级评价}}
\author{HySim-XJTU-HRPES（西安交通大学 HRPES 团队）}
\date{\today}

\begin{document}
\maketitle
\begin{abstract}
\noindent 本手册依据 HySim-XJTU-HRPES 平台场景生成模块
（\srcpath{src/scenario\_generation/}、\srcpath{include/hacdcpf/analysis/}）整理，
覆盖三类场景族——常规（负荷/新能源/储能荷电/气候扰动）、可靠性（FMEA 兼容 N-1 目录）与
弹性（台风类别驱动），并含基于 Holland 参数化风场的台风故障-修复序列生成。
注意头文件位于 \srcpath{include/hacdcpf/analysis/}（\srcpath{scenario\_generation.hpp} 等），
而非 \srcpath{include/hacdcpf/scenario\_generation/}（不存在），命名空间为 \texttt{hacdcpf::analysis}。
全部结论以源码与既有契约 \srcpath{docs/scenario\_generation\_contract.md} 为准；台风风场、修复次序、
坐标与线段为参数化近似，弹性事件通过显式标志位（\fld{used\_catalog\_sample} 等）如实声明近似来源，
候选概率仅在族内归一。文档版式与《碳流追踪与碳排放核算技术手册》一致。
\end{abstract}

\tableofcontents
\newpage

\section{阅读指南与约定}
\subsection{数据来源}
\begin{itemize}
  \item \srcpath{include/hacdcpf/analysis/scenario\_generation.hpp} —— 三类场景族入口、
        选项/结果结构与 JSON I/O；
  \item \srcpath{include/hacdcpf/analysis/typhoon\_resilience.hpp} —— Holland 风场、
        线段易损度与台风故障序列；
  \item \srcpath{include/hacdcpf/analysis/typhoon\_traffic\_impact.hpp} —— 台风对交通的
        耦合影响；
  \item \srcpath{src/scenario\_generation/scenario\_generation.cpp}、
        \srcpath{typhoon\_resilience.cpp}、\srcpath{typhoon\_traffic\_impact.cpp} —— 实现；
  \item \srcpath{docs/scenario\_generation\_contract.md} —— 输入/概率口径/近似标志契约；
  \item \srcpath{tests/test\_scenario\_generation.cpp} 等 —— 验证依据。
\end{itemize}
命名空间为 \texttt{hacdcpf::analysis}。

\subsection{单位制与索引空间}
时间以小时为步长，常规族默认 \fld{num\_steps}=8760（整年逐时序列）；台风时域
\fld{horizon\_hours} 默认 48\,h。风速单位 m/s，气压亏损单位 hPa，半径/长度单位 km，
经纬度为地理坐标。候选场景挂接 \fld{TimeSeriesData}，其内部索引空间与被扰动系统的
组件顺序一致；N-1 目录中的元件引用采用对外稳定 ID（经映射回原始组件），不使用内部
向量位置或图节点索引。

\subsection{诚实口径}
本模块不含单一“成功”布尔量。候选权重仅在族内归一
$\bigl(\sum_{c\in\mathcal{G}}\pi_c=1\bigr)$，聚类概率在聚合后归一，二者均非跨族的
无条件概率。台风相关量为参数化近似：弹性事件通过 \fld{used\_catalog\_sample}、
\fld{used\_category\_fallback}、\fld{used\_approximate\_repair\_order}、
\fld{used\_fallback\_coordinates}、\fld{used\_synthetic\_segments} 显式标注其近似来源。
可靠性族刻意仅遵循已实现的可靠性 FMEA 目录，\fld{include\_vpps}/\fld{include\_microgrids}
默认为假且兼容标志不会将其扩展到 \texttt{reliability\_assessment} 之外。可复现性仅在
系统/选项/资源/实现固定时成立。

\section{功能定位与架构}
\subsection{公共入口族}
\begin{center}\footnotesize
\begin{tabular}{@{}L{6.6cm}L{3.0cm}L{5.4cm}@{}}
\toprule
入口 & 定义位置 & 说明\\
\midrule
\srcpath{generate\_scenarios(sys, options)} & scenario\_generation.hpp & 运行所有已启用族（须至少启用一族）\\
\srcpath{generate\_regular\_scenarios(...)} & 同上 & 构建常规族（扰动时序）\\
\srcpath{generate\_reliability\_scenarios(...)} & 同上 & 构建可靠性族（每 N-1 一组）\\
\srcpath{generate\_resilience\_scenarios(...)} & 同上 & 构建弹性族（每台风类别一组）\\
\srcpath{enumerate\_n1\_contingencies(sys, options)} & 同上 & 可靠性 FMEA 兼容 N-1 目录\\
\srcpath{scenario\_generation\_options\_from\_json(j)} & 同上 & JSON 选项解析\\
\srcpath{scenario\_generation\_result\_to\_json(result)} & 同上 & 结果 JSON 序列化\\
\srcpath{generate\_typhoon\_fault\_sequence(sys, opts)} & typhoon\_resilience.hpp & Holland 轨迹→段易损→故障/修复序列\\
\srcpath{generate\_typhoon\_line\_segments(sys, opts, ...)} & 同上 & 线路自动分段（风雨空间采样）\\
\bottomrule
\end{tabular}
\end{center}

\subsection{关键数据结构}
\compmeta{ScenarioGenerationOptions}{include/hacdcpf/analysis/scenario\_generation.hpp}
嵌套 \texttt{RegularScenarioOptions}（SSP/年份、\fld{candidate\_count}、\fld{cluster\_count}、
\fld{num\_steps}=8760）、\texttt{ReliabilityScenarioOptions}（各元件包含标志、
\fld{candidates\_per\_contingency}）、\texttt{ResilienceScenarioOptions}
（\fld{intensity\_levels}、\fld{month}）、\texttt{PerturbationOptions}、
\texttt{TyphoonImpactOptions} 与 \texttt{ClusteringOptions}（\fld{method}=
\texttt{hybrid\_kmedoids\_tail\_5pct}、\fld{tail\_fraction}=0.05、
\fld{outage\_hamming\_weight}、\fld{include\_tail\_anchors}、\fld{freeze\_anchors}、\fld{seed}）。

\compmeta{ScenarioCandidate}{include/hacdcpf/analysis/scenario\_generation.hpp}
主要字段：\fld{family}、\fld{probability}、\fld{time\_series}（\texttt{TimeSeriesData}）、
可选 \fld{contingency}/\fld{resilience\_event}、\fld{outage\_signature}、\fld{risk\_score}、
\fld{tail\_anchor}。

\compmeta{TyphoonScenarioOptions}{include/hacdcpf/analysis/typhoon\_resilience.hpp}
主要字段：\fld{horizon\_hours}=48、\fld{requested\_category}（\texttt{TD..SuperTY}）、
\fld{auto\_segment\_lines}、\fld{target\_segment\_length\_km}、\fld{min\_fault\_probability}、
\fld{apply\_pv\_wind\_derating}。轨迹点 \texttt{TyphoonTrackPoint} 含
\fld{latitude}/\fld{longitude}、\fld{delta\_p\_hpa}、\fld{holland\_b}、\fld{rmw\_km}、
\fld{translation\_speed\_kmph}、\fld{vmax\_ms}、\fld{sst\_c}；段易损
\texttt{TyphoonSegmentFragility} 含对数正态风易损（\fld{lognormal\_median\_wind\_ms}、
\fld{lognormal\_sigma}）、风-雨联合致灾（\fld{seg\_a\_wind}、\fld{seg\_b\_rain}、
\fld{seg\_c\_bias}）与电缆水文（\fld{curve\_number\_cn}、\fld{manning\_n}）。

\section{核心数学模型}
场景生成在被扰动的工程系统之上构造代表性时序候选，并对高风险事件（N-1、台风）
生成故障-修复序列；AC/DC 与各元件族按契约分别处理。

\subsection{三类场景族与候选概率}
常规族在整年 \fld{num\_steps}=8760 时序上叠加负荷、新能源出力、储能荷电与气候
（SSP/年份）扰动；可靠性族为枚举出的每个 FMEA 兼容 N-1 故障建一组；弹性族为每个请求的
台风类别建一组。族内候选权重按下式归一，聚类概率在聚合后再次归一：
\begin{equation}
\pi_c = \frac{w_c}{\sum_{c'\in\mathcal{G}} w_{c'}},\qquad \sum_{c\in\mathcal{G}}\pi_c = 1 .
\end{equation}

\subsection{台风 Holland 参数化风场}
沿台风轨迹，径向气压廓线与梯度风采用 Holland(1980) 参数化：
\begin{align}
p(r) &= p_c + \Delta p \, \exp\!\Bigl[-\bigl(\tfrac{R_{mw}}{r}\bigr)^{B}\Bigr],\\
V_g(r) &= \sqrt{\frac{B\,\Delta p}{\rho_a}\Bigl(\frac{R_{mw}}{r}\Bigr)^{B}
          \exp\!\Bigl[-\bigl(\tfrac{R_{mw}}{r}\bigr)^{B}\Bigr]
          + \Bigl(\frac{r f}{2}\Bigr)^{2}} - \frac{r f}{2},
\end{align}
其中 $\Delta p$=\fld{delta\_p\_hpa}（气压亏损）、$B$=\fld{holland\_b}（形状参数）、
$R_{mw}$=\fld{rmw\_km}（最大风速半径）、$f$ 为科氏参数、$\rho_a$ 为空气密度；
移动风场按平移速度 \fld{translation\_speed\_kmph} 矢量叠加。风场按线路自动分段
（目标长度 \fld{target\_segment\_length\_km}）逐段空间采样。

\subsection{段易损度与故障-修复序列}
段级失效由对数正态风易损与风-雨联合致灾共同决定：
\begin{equation}
P_f^{\text{wind}}(v) = \Phi\!\left(\frac{\ln v - \ln \tilde v}{\sigma}\right),\qquad
P_{\text{seg}} = \frac{1}{1+\exp\!\bigl[-(a\,v + b\,R + c)\bigr]},
\end{equation}
其中 $\tilde v$=\fld{lognormal\_median\_wind\_ms}、$\sigma$=\fld{lognormal\_sigma}、
$\Phi$ 为标准正态 CDF，$v$/$R$ 为段处风速/降雨，$a,b,c$ 分别为 \fld{seg\_a\_wind}、
\fld{seg\_b\_rain}、\fld{seg\_c\_bias}；电缆通道另由曲线数 \fld{curve\_number\_cn} 与
糙率 \fld{manning\_n} 表征水文致灾。低于 \fld{min\_fault\_probability} 的段被过滤，
其余段抽样得到概率性故障并排定分阶段修复次序。

\subsection{混合 k-medoids 聚类与尾部锚定}
候选由带 5\% 尾部锚定的混合 k-medoids 削减为代表集：
\begin{equation}
\min_{\mathcal{M},\,|\mathcal{M}|=K}\ \sum_{c}\ \min_{m\in\mathcal{M}} d(c,m),
\end{equation}
距离 $d$ 融合时序距离与停运汉明距离（权重 \fld{outage\_hamming\_weight}）；
风险最高的 \fld{tail\_fraction}（0.05）候选作为尾部锚点（\fld{include\_tail\_anchors}/
\fld{freeze\_anchors}）强制保留，避免高影响低概率事件被平滑掉。

\section{接口与参数}
\begin{paramtable}
\fld{num\_steps} & $T$ & 步 & 8760 & 8760 & 常规族年时序步数\\
\fld{candidate\_count} & — & — & 整数 & — & 每族候选场景数\\
\fld{cluster\_count} & $K$ & — & 整数 & — & 聚类后保留代表场景数\\
\fld{tail\_fraction} & — & — & $0\sim1$ & $0.05$ & 尾部锚定比例\\
\fld{candidates\_per\_contingency} & — & — & 整数 & — & 每个 N-1 故障的候选数\\
\fld{intensity\_levels} & — & — & 整数 & — & 弹性族强度等级数\\
\fld{horizon\_hours} & $H$ & h & 整数 & $48$ & 台风故障序列时域\\
\fld{target\_segment\_length\_km} & — & km & 正数 & — & 线路自动分段目标长度\\
\fld{min\_fault\_probability} & — & — & $0\sim1$ & — & 段失效概率下限（过滤）\\
\fld{delta\_p\_hpa} & $\Delta p$ & hPa & 正数 & — & 台风中心气压亏损\\
\fld{holland\_b} & $B$ & — & $1\sim2.5$ & — & Holland 风场形状参数\\
\fld{rmw\_km} & $R_{mw}$ & km & 正数 & — & 最大风速半径\\
\fld{lognormal\_sigma} & $\sigma$ & — & 正数 & — & 对数正态易损度对数标准差\\
\fld{seed} & — & — & 整数 & — & 随机种子（复现用）\\
\end{paramtable}
选项经 \srcpath{scenario\_generation\_options\_from\_json} 解析、结果经
\srcpath{scenario\_generation\_result\_to\_json} 序列化；聚类方式固定为
\texttt{hybrid\_kmedoids\_tail\_5pct}。

\section{验证与边界局限}
\subsection{验证与校核}
注册测试：\srcpath{test\_scenario\_generation}、\srcpath{test\_typhoon\_traffic\_impact}；
场景 JSON/工作簿模式由 \srcpath{test\_scenario\_bundle\_schema} 校核，覆盖三族生成、
台风-交通耦合与导出模式一致性。

\subsection{边界与局限（如实标注）}
\begin{itemize}
  \item 可靠性故障刻意仅遵循已实现的可靠性 FMEA 目录；\fld{include\_vpps}/
        \fld{include\_microgrids} 默认为假，兼容标志不会将其扩展到
        \texttt{reliability\_assessment} 之外；
  \item 候选权重族内归一、聚类概率聚合后归一，均非跨族无条件概率；
  \item 可复现性仅在系统/选项/资源/实现固定时成立；
  \item 台风风场/修复次序/坐标/线段为参数化近似，弹性事件以
        \fld{used\_catalog\_sample}、\fld{used\_category\_fallback}、
        \fld{used\_approximate\_repair\_order}、\fld{used\_fallback\_coordinates}、
        \fld{used\_synthetic\_segments} 如实声明近似来源。
\end{itemize}
\section{跨模块工业级技术评价基线}

本章规定本手册所述模块进入工程算例、批量研究或生产服务前必须共同满足的评价基线。
具体模型公式、变量、约束和专用算法以正文相应章节为准；本章不扩大源码已经实现的范围。

\subsection{输入、输出、边界条件与数据结构审计}
输入审计应逐项检查有限性、单位、上下界、枚举值和引用完整性。系统对象必须先按所需层级完成
校验；AC 与 DC 母线采用域限定映射，组件稳定 \texttt{index}、向量位置和临时图索引不得混用。
输出审计必须记录单位、索引空间、模型范围、收敛或可行状态、后端、容差、迭代/节点计数和告警。
空系统、空时域、孤岛、重复 ID、零容量、退化约束与不可用可选后端均属于边界测试集。

\subsection{工程假设与数值稳定性分析}
模型使用的平衡、线性化、稳态、独立性、概率分布、时间离散或网络等值假设必须与应用场景匹配。
数值评价至少包含变量缩放、绝对/相对残差、可行性违反、条件数或枢轴质量、步长/阻尼、迭代停滞
和随机种子复现性。若采用正则化、平滑、回退、松弛或近似，应同时报告触发条件、参数值、对主指标
的敏感性以及是否改变模型口径；不得用“已返回结果”替代收敛或有效性证明。

\subsection{复杂度与资源分析}
令 $n_x$、$n_c$、$n_t$、$n_s$、$|V|$、$|E|$ 分别表示变量、约束、时段、场景、节点和边数。
报告应分别给出建模、求解、后处理的时间与内存。图遍历通常为 $O(|V|+|E|)$；逐时段/逐场景
流程至少线性增长为 $O(n_t n_s)$ 次子问题调用；稠密线性代数最坏可达 $O(n_x^3)$，稀疏方法则应
报告结构非零数、填充率和因子分解峰值内存；MILP/NLP 不给出虚假的多项式最坏界，而报告节点数、
间隙、时限和实测扩展曲线。复杂度结论必须基于固定硬件、固定后端和可复现算例族。

\subsection{异常工况处理}
非法输入应在进入求解器前返回结构化错误；不可行、无界、不收敛、时限、内存不足和后端缺失必须
相互区分。部分结果只有在对应 \fld{ValidityFlags}、\fld{model\_scope}、
\fld{model\_limitations} 或等价字段允许时才可使用。异常路径不得静默丢弃 AC/DC 资产、制造平衡
节点、把 NaN 改写为零或把近似解标为全局最优；系统原始对象应保持可恢复和可重试。

\subsection{与其他模块的接口关系}
接口链路遵循“工程语义模型 $\to$ 校验 $\to$ 投影/装配 $\to$ 求解或分析 $\to$ 结果回投影”。
跨模块传递时保留稳定 ID、单位、符号、时间基准和场景身份；消费潮流、OPF、动态或可靠性结果前，
先检查其收敛、覆盖范围和有效性标志。HTTP/JSON/Python 层只做语义保持适配，不重新解释求解结果。

\subsection{测试与验证建议}
最低验证矩阵包括：手算小例、标准算例、边界/异常输入、随机性质测试、算法或后端交叉校核、
序列化往返、稳定 ID 回投影、AC/DC 同号母线、重复运行确定性和规模扩展测试。数值算法还应加入
残差独立重算、雅可比有限差分或自动微分校核；近似模型应与高保真模型比较并给出误差分布，而非
只报告单个平均值。回归报告必须列明构建配置、算例版本、容差、通过/失败/跳过数及未验证范围。

\subsection{工业级评估指标}
\begin{itemize}
  \item \textbf{正确性}：守恒残差、约束最大违反、解析/基准误差、结果回投影一致性；
  \item \textbf{鲁棒性}：收敛率、不可行识别率、异常分类准确率、扰动敏感性与随机复现率；
  \item \textbf{性能}：端到端时延、吞吐量、峰值内存、并行效率与规模增长斜率；
  \item \textbf{可追溯性}：输入模式版本、稳定 ID 覆盖率、模型范围/限制字段完整率、告警可定位率；
  \item \textbf{工程有效性}：与标准、外部引擎或现场数据的偏差，以及误判/漏判对业务指标的影响。
\end{itemize}
指标应预先给出验收阈值和失败处置；未达到阈值时只能标记为实验性或受限适用。

\subsection{适用范围、局限性与后续改进}
本手册只评价当前源码覆盖的模型、后端和数据结构，不对未建模物理过程、未安装求解器或未执行的
外部验证作保证。后续改进应按“缺口 $\to$ 可量化指标 $\to$ 源码/测试变更 $\to$ 独立验证”闭环，
优先消除会造成错误结论的语义缺口，其次提升数值鲁棒性与性能，最后扩展模型范围。


\end{document}
